% Nouveau EGA V (cote 25): archive cote 25, site page 39 (PDF page 40).
% typescript. Draft compared: none.
% AI-transcribed by Claude Opus 5.5 (claude-opus-5-5, adaptive thinking, effort high), Message Batches API, batch ega25b (msgbatch_01A7n5YuXubUPNK1zsepi8RG), 2026-11-13; not yet checked by eye.
\scanpage{25/39}
\mspage{\hand{\unsure{35|33}}}

une sous-variété linéaire passant par a son espace tangent en a, consi-%
déré comme sous-espace \xout{tangent} de l'espace tangent à P en a. Ce genre
\margin{\hand{|}}de \unsure{sorites|soritites}, ensemble avec divers sorites sur les sous-variétés linéaires
et les grassmaniennes, devraient être donnés dans un ou deux n$^{\text{o}}$s prélimi
naires, bien entendu en énonçant tout systématiquement sur une base
quelconque.

En fait, il y a mieux, savoir que le \uline{sous-préschéma} $\underline{Y}^{\mathrm{sing}}$
des zéros singuliers de $\phi$ relativement à $\overset{\hand{$\vee$}}{P}$ , défini dans par.16, n'est
autre que $P(\hand{$\unsure{\underline{\nu}}$}_{X/P})$ , et (comme ce dernier est lisse sur S de dimension
relative d \unsure{$+$} (r-d-1) = r - 1 (r étant la dim relative de P sur S)
on est sous les conditions favorables étudiées au par.16 . \deleted{Pour vérifi
er cette assertion, on utilise le fait suivant (qui devrait figurer au
par. 17 des morphismes lisses), en l'appliquant à $X_{\check{P}}$ \inserted[typed]{et $\underline{P}^1(\underline{O}\ (1,1))$} S : Soit
X lisse sur S de dimension relative m, $\underline{M}$ un Module localement libre sur
X de rang n, $\mu$ une section de $\underline{M}$ , x V( $\mu$ )=Z , supposons que
dim$_x$ Z$_s$ = n-m , \xout{et} (resp. et que $Z_{\underline{s}\ \text{réd}}$ soit lisse sur $\overline{k(s)}$ en $\bar{x}$ )
alors Z est \xout{lisse sur} plat sur S en x (resp. lisse sur S en x). (Dém:
on peut supposer M = $\underline{O}_X^{\,m}$ , donc $\mu = (\mu_1,..,\mu_n)$ , alors il résulte de
la théorie de la dimension et du fait que $\underline{O}_{X_s,x}$ est CM que les
$\mu_{i\ s}$ forment une suite $\underline{O}_{X_s}$-régulière en x, donc \deleted{les $\mu_i$ sont transvers
salement $\underline{O}_X$-réguliers rel. à S; di de plus Z}} Pour le vérifier,
notons que par définition, $\underline{Y}^{\mathrm{sing}}$ n'est autre que le sous-préschéma \inserted{de $\underline{Y}$}
des zéros de la section \inserted{$\psi =$} $d\phi|\underline{Y}$ de $\underline{\Omega}^1_{X_{\check{P}}/\check{P}}(1,1)\otimes \underline{O}_Y = \underline{\Omega}^1_{X/S}\otimes\underline{O}_Y(1,1)$\hand{.}
Nous allons donner une autre interprétation de cette section, dont
la conclusion résultera aussitôt. Pour ceci, considérons le diagramme
\deleted{de Modules} \inserted{\unsure{suites exactes}}
suivant sur $X_{\check{P}}$ , ou plus généralement sur un préschéma Z \inserted{quelconque} au dessus de
$X_{\check{P}}$ :

\mspage{\hand{39}}
